\subsection{Limitations and Future Work}

The present study has three main limitations that define the scope of the reported conclusions. First, the available real-flight supervision remains limited to telemetry anomalies and fault-domain labels. UAV-SEAD does not provide source-level software fault annotations, while the ALFA evaluation is restricted to binary anomaly transfer and the Uncategorized flights lack temporal event labels. In addition, the decrease in Stage 1 performance under flight-level isolation indicates that anomaly screening remains sensitive to cross-flight distribution shifts. Future work should therefore prioritize stronger cross-flight anomaly models and evaluation on additional independently collected flight logs, while retaining complete flight-level separation for model selection and testing.

Second, the architecture-aware evidence is based on static and associative relationships. The current telemetry-to-subsystem mapping is manually defined from documented signal semantics, and the PX4 module graph is derived from commit-specific uORB publisher--subscriber relations rather than runtime execution provenance. This representation provides a reproducible bridge from telemetry evidence to source-tree inspection candidates, but it cannot determine whether a ranked module was causally responsible for the observed anomaly. A stronger validation framework would combine independently adjudicated semantic mappings with build-specific dependency resolution and runtime topic provenance, allowing static architectural evidence to be compared with the modules that are actually active during a failure.

Third, the matched-reference replay extension is evaluated under a constrained controlled setting. It requires an additional healthy replay of the same source flight, covers four mutation mechanisms at a single PX4 revision, and is based on a limited number of independent development and transfer source flights. Moreover, the current residual construction and evidence aggregation operate offline, and the existing transfer study does not constitute a fully untouched end-to-end generalization test because some preprocessing components had previously encountered the transfer sources. Future replay experiments should use completely independent calibration and evaluation flights, expand the number of source logs, firmware revisions, and mutation types, and freeze the full preprocessing and calibration pipeline before evaluation. Replacing full-log support checks with causal online alignment would also be necessary before the replay extension could be assessed for real-time use.

These limitations do not affect the interpretation of the recorded-data classification and attribution results within their stated evaluation protocols, but they define the steps required to move from telemetry-driven diagnostic triage toward stronger software-level fault localization. The most important next directions are therefore to improve cross-flight anomaly screening, replace static architectural associations with runtime-aware evidence, and validate module rankings on broader and fully independent source-level fault scenarios.
